% ==============================================================
\section{結論}\label{sec:conclusion}
% ==============================================================

本論文では、Apriori-Window~\cite{aprioriwindow}を
任意の$d$次元グリッドに拡張する
\textbf{Multi-Dimensional Apriori-Window}を提案した。
固定空間始点ごとの1次元マルチセット還元、
重複考慮$t$軸ストライド、
およびシングルトン区間マップ交差枝刈りの3技法により、
全密集パターンとその任意形状の密集領域を厳密に列挙する。

実験により以下を示した：
(1)~4つの実データセット（$d{=}3$--$4$）において、
既存手法はパターン検出F1 $\leq$ 0.71、
領域復元平均Jaccard $\leq$ 0.104にとどまった；
(2)~合成データにおいて$N = 10^3$--$10^8$で
線形スケーリングを確認し、
$10^8$トランザクションを10分以内に処理可能であった。

\paragraph{制約と限界.}
(1)~本手法はグリッド離散化を前提としており、
ビン幅および$\bm{W}$の選択が結果に影響する。
(2)~$W_{\mathrm{sp}}$や$d$が増加すると
マルチセット構築の$O(N \cdot W_{\mathrm{sp}})$コストが増大し、
高次元設定では次元削減や粗いビン化が必要となる。
(3)~本手法は検出と記述に特化しており、
因果的説明や外部イベントとの関連付けは対象外である。

\paragraph{今後の課題.}
(1)~\emph{適応的パラメータ選択}：
現在、ウィンドウサイズ$\bm{W}$と閾値$\theta$は
ユーザーが事前に指定する。
局所的なデータ統計量（密度分布、空間的自己相関）に基づき
$\bm{W}$と$\theta$を自動調整するフレームワークは
実用上重要である。
(2)~\emph{更なる高速化}：
空間始点ごとの時間マルチセット構築が
主要な計算ボトルネックである。
累積和に基づく$O(1)$ウィンドウ集約や、
空間軸に沿った差分更新（スライス差分）による
冗長計算の排除が有望である。
(3)~\emph{近似手法}：
厳密解が不要な探索的分析では、
空間始点をランダムサンプリングして
走査対象を間引くことで、
$W_{\mathrm{sp}}$が大きい高次元設定での計算コストを削減できる。
サンプリング率と見落とし率のトレードオフの
理論的評価は今後の課題である。
